\documentclass[11pt,a4paper]{article}
\usepackage[margin=18mm]{geometry}
\usepackage{amsmath,amssymb,mathtools}
\usepackage{booktabs,array,enumitem}
\usepackage[dvipsnames]{xcolor}
\usepackage{listings}
\usepackage{tikz}
\usetikzlibrary{arrows.meta,positioning}
\usepackage[most]{tcolorbox}
\usepackage{hyperref}
\hypersetup{colorlinks=true,urlcolor=MidnightBlue,linkcolor=MidnightBlue}
\setlength{\parindent}{0pt}
\setlength{\parskip}{4pt}
\newcommand{\course}{Computational Methods in Physics}
\newtcolorbox{keybox}{colback=RoyalBlue!5,colframe=RoyalBlue!65!black,boxrule=.7pt,arc=2mm}
\newtcolorbox{checkbox}{colback=ForestGreen!5,colframe=ForestGreen!55!black,boxrule=.7pt,arc=2mm}
\newtcolorbox{aibox}{colback=BurntOrange!7,colframe=BurntOrange!80!black,boxrule=.7pt,arc=2mm}
\lstdefinestyle{matlabstyle}{
  language=Matlab,
  basicstyle=\ttfamily\footnotesize,
  keywordstyle=\color{RoyalBlue!75!black},
  commentstyle=\color{ForestGreen!55!black},
  stringstyle=\color{BurntOrange!80!black},
  numbers=left,
  numberstyle=\scriptsize\color{black!50},
  numbersep=5pt,
  frame=single,
  rulecolor=\color{RoyalBlue!25},
  backgroundcolor=\color{RoyalBlue!3},
  showstringspaces=false,
  columns=fullflexible,
  keepspaces=true,
  breaklines=true
}
\begin{document}

{\large\textsc{\course}}\\[-2pt]
{\small Learning Note \textbar\ Week 5}\par\vspace{4pt}
{\LARGE\bfseries Root Finding: Bracketing, Newton, and Reliability}\par\vspace{3pt}
\textit{A root is not trustworthy because an iteration stopped; it is trustworthy when the model, convergence evidence, residual, and physical interpretation agree.}

\begin{keybox}
\textbf{Physical question.} For an elliptic orbit, how can we solve Kepler's nonlinear equation reliably, and what evidence tells us that a computed eccentric anomaly is genuinely a root?
\end{keybox}

\section*{Learning outcomes}
By the end of this week, you should be able to:
\begin{enumerate}[nosep,leftmargin=5mm]
  \item explain why nonlinear equations usually require iterative root finding;
  \item construct and justify a sign-changing bracket, then apply bisection with a convergence guarantee;
  \item derive and apply Newton's update while analysing initial-guess and derivative sensitivity;
  \item distinguish step size, residual, root-location error, iteration limit, and solver status;
  \item diagnose invalid brackets, small derivatives, divergence, cycling, and false convergence;
  \item use MATLAB \texttt{fzero} appropriately and independently validate its result; and
  \item compare methods quantitatively using iteration count, residual, robustness, and available physical evidence.
\end{enumerate}

\section*{1. Turn a nonlinear physical equation into a residual}
Many physical models ask for a quantity that is defined implicitly rather than by a formula of the form $x=g(\hbox{known data})$. Equilibria, threshold conditions, event times, calibration equations, and transcendental relations often reduce to
\begin{equation}
 f(x)=0.
\end{equation}
The function $f$ is a \textbf{residual}: it measures how far a trial value is from satisfying the defining equation. A nonlinear residual may contain products, powers, exponentials, logarithms, or trigonometric functions, so ordinary linear-system methods do not isolate the answer in one algebraic step.

For an elliptic orbit, Kepler's equation connects mean anomaly $M$, eccentric anomaly $E$, and eccentricity $e$:
\begin{equation}
 M=E-e\sin E,\qquad 0\le e<1.
\end{equation}
Angles $E$ and $M$ are measured in radians and $e$ is dimensionless. Define
\begin{equation}
 f(E)=E-e\sin E-M,
 \qquad
 f'(E)=1-e\cos E.
\end{equation}
The computational problem is now explicit: find $E$ for which $f(E)=0$.

\begin{checkbox}
\textbf{Predict before computing.} If a numerical method returns an $E$ that changes by only $10^{-12}$ rad between two iterations, is that alone enough evidence that Kepler's equation is satisfied? What else would you inspect?
\end{checkbox}

\section*{2. Bracketing gives global existence evidence}
Suppose $f$ is continuous on $[a,b]$ and $f(a)f(b)<0$. The Intermediate Value Theorem guarantees at least one root between the endpoints. This is a \textbf{bracket}: a global invariant that does not depend on a tangent or a visually good guess.

Kepler's equation provides an unusually useful physical bracket. Choose
\begin{equation}
 a=M-e,\qquad b=M+e.
\end{equation}
Then
\begin{align}
 f(M-e)&=-e-e\sin(M-e)=-e[1+\sin(M-e)]\le0,\\
 f(M+e)&= e-e\sin(M+e)= e[1-\sin(M+e)]\ge0.
\end{align}
Because $f$ is continuous, a root is guaranteed in $[M-e,M+e]$. We can also establish \textbf{uniqueness}. Since $\cos E\le1$,
\begin{equation}
 f'(E)=1-e\cos E\ge 1-e>0 \qquad (e<1).
\end{equation}
Thus $f$ is strictly increasing: the bracket contains exactly one root. Notice the division of evidence: the sign change establishes existence, while monotonicity establishes uniqueness.

\section*{3. Bisection trades speed for a guarantee}
Bisection repeatedly halves a valid bracket. At iteration $k$, form
\begin{equation}
 c_k=\frac{a_k+b_k}{2},
\end{equation}
evaluate $f(c_k)$, and keep the half-interval whose endpoints still have opposite signs. The root never leaves the retained interval.

If the current bracket width is $w_k=b_k-a_k$, the midpoint has a certified root-location bound
\begin{equation}
 |c_k-E_*|\le\frac{w_k}{2},
\end{equation}
where $E_*$ is the unknown root. After each successful sign update the width halves, so the method converges predictably even when the residual curve is awkward.

\begin{lstlisting}[style=matlabstyle,caption={Core bisection logic for one iteration.}]
c = a + 0.5*(b-a);
fc = f(c);
if f(a)*fc <= 0
    b = c;
else
    a = c;
end
\end{lstlisting}

\begin{keybox}
\textbf{Do not run bisection on an invalid bracket.} If the endpoint residuals have the same sign, the bisection guarantee is absent. More iterations do not repair missing existence evidence; first change or justify the interval.
\end{keybox}

\section*{4. Stopping criteria answer different questions}
An iterative solver needs a deliberate stopping rule. Common evidence includes:
\begin{center}
\begin{tabular}{>{\bfseries}p{31mm}p{54mm}p{76mm}}
\toprule
Evidence & Typical quantity & What it tells you \\
\midrule
Bracket width & $(b-a)/2$ & A certified root-location bound for a valid bisection bracket. \\
Step size & $|x_{k+1}-x_k|$ & The iteration has stopped moving much; not automatically a small residual. \\
Residual & $|f(x_k)|$ & How well the defining equation is satisfied; relation to root error depends on local slope. \\
Iteration limit & $k\le k_{\max}$ & Prevents an infinite or excessively expensive run; reaching it is a failure signal, not convergence. \\
Solver status & exit flag/message & Reports what the implementation detected; still requires model and residual checks. \\
\bottomrule
\end{tabular}
\end{center}
A robust Newton stopping decision can require both
\begin{equation}
 |x_{k+1}-x_k|\le \tau_x(1+|x_{k+1}|)
 \quad\text{and}\quad
 |f(x_{k+1})|\le\tau_f.
\end{equation}
The first mixes absolute and relative scale; the second directly tests the equation. Tolerances should ultimately reflect the physical accuracy required, not merely the smallest number the computer can print.

\section*{5. Newton's method is fast local reasoning}
Near a current estimate $x_k$, approximate the residual by its tangent:
\begin{equation}
 f(x)\approx f(x_k)+f'(x_k)(x-x_k).
\end{equation}
Set the tangent approximation to zero and solve for the next estimate:
\begin{equation}
 \boxed{x_{k+1}=x_k-\frac{f(x_k)}{f'(x_k)}}.
\end{equation}
For a smooth simple root and a sufficiently good initial guess, Newton often converges far faster than bisection. The price is that it relies on \textbf{local} information. If $f'(x_k)$ is tiny, the division can create an enormous step. A poor initial guess can also send the iteration away from the intended root, into a cycle, or into a region where the model/function is invalid.

For Kepler's equation, a natural baseline initial guess is $E_0=M$. With $e=0.70$ and $M=1.00$ rad, Newton rapidly approaches
\begin{equation}
 E_*\approx1.6946389\ \mathrm{rad}.
\end{equation}
Bisection reaches the same root more slowly but carries a valid bracket throughout.

\section*{6. A stress case exposes derivative sensitivity}
Take a highly eccentric orbit with
\begin{equation}
 e=0.999,\qquad M=0.15\ \mathrm{rad},\qquad E_0=0.
\end{equation}
At the initial point,
\begin{equation}
 f(0)=-0.15,\qquad f'(0)=1-0.999=0.001.
\end{equation}
The first raw Newton step is therefore
\begin{equation}
 \Delta E=-\frac{-0.15}{0.001}=150\ \mathrm{rad}.
\end{equation}
That huge step is a clear warning: the local tangent is poorly conditioned for this starting point. It does \emph{not} prove that raw Newton can never recover, but it proves that ``Newton is always safe because it is fast'' is indefensible.

One remedy is \textbf{safeguarding}: retain a valid bracket, propose a Newton step, and reject it if the candidate leaves the bracket or the derivative is unusable. A midpoint step then restores progress. This combines Newton's fast local information with the global reliability of bracketing.

\begin{aibox}
\textbf{Responsible use of AI-generated numerical code.} If an AI suggests a root finder, do not validate it only by obtaining a plausible number. Test the bracket logic, derivative formula, stopping condition, iteration limit, final residual, and at least one independent reference or limiting case. Record which material suggestion you accepted, modified, or rejected and why.
\end{aibox}

\section*{7. MATLAB fzero is a reference implementation, not an oracle}
MATLAB's base function \texttt{fzero} is the appropriate official scalar root finder for many production calculations. When a sign-changing interval is available, pass the bracket rather than discarding that physical evidence:
\begin{lstlisting}[style=matlabstyle,caption={Reference solve with an independently checked residual.}]
f = @(E) E - e*sin(E) - M;
bracket = [M-e, M+e];
[Eroot,fval,exitflag,output] = fzero(f,bracket);
residual = abs(f(Eroot));
\end{lstlisting}
Use the explicit Week 5 bisection and Newton implementations to understand algorithm behaviour; use \texttt{fzero} to compare against a tested MATLAB implementation. Agreement among independent methods is useful evidence, but they can all agree on the wrong physical model, so assumptions and units still matter.

\section*{8. Compare methods quantitatively}
For the baseline Kepler problem, a useful comparison table contains more than the root:
\begin{center}
\begin{tabular}{p{25mm}p{30mm}p{32mm}p{62mm}}
\toprule
Method & Main input evidence & Useful output evidence & Main reliability issue \\
\midrule
Bisection & Valid sign-changing bracket & Iterations, residual, final half-width & Guaranteed but comparatively slow. \\
Newton & Initial guess and derivative & Iterations, step, residual, derivative behaviour & Can take dangerous steps or fail from poor local information. \\
\texttt{fzero} & Preferably a bracket & Exit flag, function count, residual & Tested implementation still cannot validate the physical model for you. \\
\bottomrule
\end{tabular}
\end{center}
Iteration count measures efficiency, not correctness. Residual measures equation satisfaction, not automatically root error. Robustness requires controlled tests: change the bracket, initial guess, eccentricity, or tolerance and inspect whether the conclusion remains credible.

\section*{9. Practical transfer and capstone progression}
In the practical, you will compare explicit bisection, explicit Newton, safeguarded Newton, and \texttt{fzero}; diagnose an invalid bracket; and vary eccentricity while tracking the derivative bound $1-e$. Preserve a compact evidence record suitable for the Weeks 01--06 mid-semester report.

For your capstone question, look for a scalar condition of the form $f(x)=0$: an equilibrium, threshold, event time, calibration parameter, resonance condition, or transcendental relationship. Before the Week 06 feasibility proposal, record (i) what the root means physically, (ii) whether a bracket or justified initial guess is available, (iii) the intended stopping criterion, and (iv) at least two independent validation checks.

\section*{Three takeaways}
\begin{keybox}
\begin{enumerate}[nosep,leftmargin=5mm]
  \item \textbf{Bracket first when you can.} A verified sign change gives global existence evidence; bisection converts it into a certified geometric uncertainty.
  \item \textbf{Newton is fast because it trusts local geometry.} That strength becomes a weakness when the initial guess or derivative is poor.
  \item \textbf{Stopping is an evidence decision.} Keep step size, residual, root error, iteration limit, solver status, and physical plausibility conceptually separate, then combine the evidence appropriate to the problem.
\end{enumerate}
\end{keybox}

\begin{keybox}
\textbf{Preparation for the practical.} Bring one clear prediction about when Newton should outperform bisection and one situation in which you would prefer a bracketed method. Be ready to justify your choice using available physical evidence, a stopping rule, and an independent validation check rather than method reputation alone.
\end{keybox}

\end{document}
